\chapter{Computational Quantities: Units, Coordinates, and Metadata}
\index{computational quantity|textbf}
\index{unit!computational quantity|textbf}
\index{metadata|textbf}

\label{ch:training-computational-quantities}

\chapterstatus{pedagogical foundation. The examples use illustrative synthetic
values and establish no material result, numerical convergence claim, scientific
validation, or uncertainty quantification.}

\section{Learning goals}

After working through this chapter, a reader should be able to

\begin{enumerate}
  \item distinguish a mathematical number from a computational quantity;
  \item attach units, coordinate meaning, indexing conventions, and provenance
        to numerical values;
  \item explain why shape and data type do not determine scientific meaning;
  \item convert units and coordinates without silently changing the represented
        quantity; and
  \item identify defining metadata that must accompany a retained array.
\end{enumerate}

This monograph begins with data rather than with a Hamiltonian. Every later
vector, matrix, grid, and operator is stored as numbers, and those numbers are
usable only when their interpretation is explicit.

\section{From numbers to quantities}
\index{number!quantity distinction}

A floating-point value such as \(5.43\) is a number. It becomes a computational
quantity only after its role and unit are supplied. A lattice length might be
recorded schematically as
\begin{equation}
  q=(5.43,\ \text{angstrom}).
\end{equation}
A conversion changes the numerical value and the unit together while preserving
the physical quantity. Changing only one member of the pair changes the
meaning.

The same rule applies to arrays. An array does not announce whether it contains
positions, reciprocal vectors, energies, expansion coefficients, or matrix
elements. Its shape and data type are necessary software facts, but they are not
a scientific interpretation.

\begin{definition}[Computational quantity]
A computational quantity is a numerical value or array together with the
metadata required to interpret it for a declared operation. Depending on the
operation, that metadata can include units, coordinate system, axis meanings,
ordering, normalization, reference value, geometry, and provenance.
\end{definition}

The phrase ``required to interpret'' is operation dependent. Adding two
energies requires compatible energy units and references. Taking a dot product
of Cartesian displacement vectors requires compatible coordinates and length
units. Comparing represented Hamiltonians requires substantially more, as later
chapters show.

\section{Dimensions, units, and reference values}
\index{dimension!physical}
\index{energy reference}

A physical dimension describes the kind of quantity, such as length, inverse
length, or energy. A unit selects a numerical scale within that dimension.
Angstrom and bohr are different length units; electronvolt and hartree are
different energy units. Unit compatibility does not imply equality of numerical
values, so conversion must be explicit.

Some quantities also depend on a reference. If an energy is reported relative
to a chosen zero \(E_0\), then
\begin{equation}
  \overline E=E-E_0.
\end{equation}
Two arrays expressed in electronvolts can still be incompatible when their
energy zeros differ. A unit conversion cannot repair a reference mismatch.

A useful dimensional audit asks three questions before arithmetic:

\begin{enumerate}
  \item Do the operands represent the same physical dimension?
  \item Are their units identical, or is an explicit conversion required?
  \item Does either quantity depend on an origin, zero, normalization, or other
        reference convention?
\end{enumerate}

Dimensionless does not mean convention free. Fractional coordinates, reduced
Bloch twists, phases, and normalized coefficients can all be dimensionless
while still requiring a basis, interval, or phase convention.

\section{Coordinates are part of the data}
\index{coordinate system!metadata}

The tuple \((1,0)\) can denote a Cartesian vector, fractional lattice
coordinates, grid indices, or coefficients in an abstract basis. These are not
interchangeable meanings. A coordinate record should identify at least

\begin{itemize}
  \item the space in which the coordinates live;
  \item the ordered axes or basis vectors;
  \item whether components are Cartesian, fractional, reciprocal, or abstract;
  \item the component unit, if any;
  \item the origin or reference point when relevant; and
  \item the vector-axis and array-order convention.
\end{itemize}

A coordinate transformation changes the numerical components while preserving
the represented object only when the transformation rule is applied
consistently. If \(A\) stores basis vectors as columns and \(\mathbf s\) stores
fractional coordinates, then
\begin{equation}
  \mathbf r=A\mathbf s.
\end{equation}
Transposing \(A\) without changing the convention is not a harmless formatting
operation. Chapter~\ref{ch:training-periodic-geometry} develops this map for
periodic cells after the required linear algebra has been introduced.

\section{Indexing, axis meaning, and ordering}
\index{array ordering}
\index{indexing convention}

Computational arrays have both mathematical axes and storage conventions. For a
table \(X\in\mathbb R^{N\times d}\), rows might enumerate points and columns
might enumerate Cartesian components. Another program may store the transpose.
Both layouts can be valid, but an unlabeled shape does not distinguish them.

Index origins also matter. Mathematical derivations often number basis vectors
from \(1\) to \(N\), whereas Python arrays use indices from \(0\) to \(N-1\).
A translation between these conventions should occur at a named boundary, not
implicitly inside a scientific argument.

For tensor-product grids, a multi-index must also be mapped to a linear index.
C ordering makes the final axis vary fastest; Fortran ordering makes the first
axis vary fastest. Two matrices assembled with different flattening conventions
can have identical shapes and spectra while assigning entries to different
physical grid points.

\section{Metadata needed by represented scientific data}
\index{represented data}

The interpretation of a finite represented operator can depend on

\begin{itemize}
  \item the identity and dimension of the represented state space;
  \item basis identity, orbital convention, and ordering;
  \item geometry, lattice vectors, site labels, and boundary convention;
  \item units and normalization;
  \item scalar energy reference;
  \item spin representation;
  \item gauge, phase, or other coordinate conventions; and
  \item provenance linking the record to its inputs and construction.
\end{itemize}

Not every elementary array needs every field. The rule is instead that the
record must carry enough information for its promised operations to fail closed.
If subtraction requires a common energy reference, an absent reference is not
permission to assume one. If a decoder promises a column-vector convention, a
row-vector document is not an alternative spelling of the same record.

It is useful to distinguish defining data from derived data. A primitive basis
can define a metric, reciprocal basis, and projector. Serializing all of them as
independent authorities creates the possibility of internal disagreement.
Retaining the defining data and deterministically recomputing derived data often
produces a stronger record.

\section{A small transparent record}

The following example keeps an energy array, its unit, and its reference in one
visible record. It uses a plain dictionary only for pedagogy; maintained project
interfaces use typed immutable records and explicit validation.

\begin{pythonexample}{converting an energy quantity without losing its reference}
import numpy as np

# These synthetic levels are measured relative to one declared energy zero.
levels_ev = np.array([0.0, 0.12, 0.31], dtype=np.float64)
record_ev = {
    "values": levels_ev,
    "unit": "eV",
    "energy_reference": "illustrative band edge",
    "axis_meaning": "state index in ascending eigensystem order",
}

# Convert the values and unit together; preserve the non-numeric metadata.
ev_per_hartree = 27.211386245988
record_hartree = {
    **record_ev,
    "values": record_ev["values"] / ev_per_hartree,
    "unit": "hartree",
}

# Recover the original numerical representation as an independent check.
np.testing.assert_allclose(
    record_hartree["values"] * ev_per_hartree,
    record_ev["values"],
)
assert record_hartree["energy_reference"] == record_ev["energy_reference"]
\end{pythonexample}

The example verifies one conversion identity for synthetic values. It does not
establish a project serialization contract or any physical energy level.

\documentcallout{Notebook to do}{Create
\path{examples/tutorials/research-monograph/foundations/computational_quantities.ipynb}
with explicit imports, explanatory inline comments, unit conversion, coordinate
annotation, malformed-metadata cases, and no hidden state. Keep it unexecuted in
the repository and verify it through a separate bounded test route.}

\section{Compatibility before arithmetic}
\index{compatibility!before arithmetic}

A robust computational sequence is
\begin{equation}
\begin{aligned}
  \text{decode}
  &\longrightarrow \text{validate each record}
  \longrightarrow \text{check pairwise compatibility}\\
  &\longrightarrow \text{perform arithmetic}
  \longrightarrow \text{retain the result context}.
\end{aligned}
\end{equation}
Validation of each operand and compatibility between operands answer different
questions. Two individually valid energy arrays can still use different units,
references, bases, or state orderings.

Automatic inference should be limited to conventions owned by a documented
contract. Inferring a unit from the magnitude of an array, an axis meaning from
its length, or a physical correspondence from matching labels replaces an
explicit scientific decision with a heuristic.

\section{Exercises}

\begin{enumerate}
  \item List the metadata needed to add two Cartesian force vectors and explain
        which fields are unnecessary for that operation.
  \item Give two different scientific interpretations of a complex
        \(4\times4\) array that have the same shape and data type.
  \item Describe why two energy arrays with identical units but different
        scalar zeros cannot be subtracted as a physical perturbation without
        alignment.
  \item Write a small record for three fractional coordinates stored by rows.
        Include dimension, axis meaning, origin convention, and provenance.
  \item Classify the direct basis, Gram metric, and reciprocal basis as defining
        or derived data under the convention developed in
        Chapter~\ref{ch:training-periodic-geometry}.
\end{enumerate}

\section{Evidence boundary}

Metadata completeness is relative to a declared operation and contract. A
successful decode or round trip establishes bounded software behavior; it does
not prove that the recorded metadata are truthful, that the source calculation
was adequate, or that two physical systems should be compared. The next two
chapters add finite coordinate representations and then Hermitian eigensystems
while preserving this distinction between numbers and represented meaning.
